% Generated draft: requires mathematical and source review.
\documentclass[11pt]{article}
\usepackage[margin=25mm]{geometry}
\usepackage{fontspec}
\setmainfont{Latin Modern Roman}
\usepackage{amsmath,amssymb,mathtools,unicode-math}
\setmathfont{Latin Modern Math}
\usepackage{enumitem,needspace,xurl,hyperref}
\setlength{\parindent}{0pt}
\setlength{\parskip}{5pt}
\setlength{\emergencystretch}{3em}
\newcommand{\N}{\mathbb{N}}
\newcommand{\Z}{\mathbb{Z}}
\newcommand{\Q}{\mathbb{Q}}
\newcommand{\R}{\mathbb{R}}
\newcommand{\C}{\mathbb{C}}
\newcommand{\F}{\mathbb{F}}
\newcommand{\A}{\mathbb{A}}
\newcommand{\PP}{\mathbb{P}}
\newcommand{\E}{\mathbb{E}}
\newcommand{\eps}{\varepsilon}
\newcommand{\dd}{\,\mathrm{d}}
\begin{document}
\section*{UnsolvedMath ID: 30002261}
\textbf{Source problem code:} catalogue entry
\subsection{Definitions and mathematical statement}
Let \(\mathcal C_{\exp}=(\C,+,\cdot,0,1,\exp)\), where \(\exp\colon\C\to\C\) is the usual complex exponential function. An uncountable first-order structure \(M\) is \emph{quasiminimal} if every subset of its underlying set definable by a first-order formula with finitely many parameters from \(M\) is countable or has countable complement. Here countable includes finite.

\begin{enumerate}
\item \textbf{Quasiminimal complex exponentiation (attributed to Zilber).} Is \(\mathcal C_{\exp}\) quasiminimal? Explicitly, for every integer \(n\geq 0\), every first-order formula \(\varphi(x,y_1,\ldots,y_n)\) in the language \(\{+,\cdot,0,1,\exp\}\), and every tuple \(\bar a\in\C^n\), set
\[
X_{\varphi,\bar a}=\{z\in\C:\mathcal C_{\exp}\models\varphi(z,\bar a)\}.
\]
Is it always the case that
\[
|X_{\varphi,\bar a}|\leq\aleph_0
\quad\text{or}\quad
|\C\setminus X_{\varphi,\bar a}|\leq\aleph_0?
\]
\item \textbf{Abelianness of regular groups.} Is every regular group abelian? For a group \(G\), being abelian means
\[
\forall g,h\in G,\qquad gh=hg.
\]
The supplied record does not define \emph{regular group}. If \(\mathcal R\) denotes the intended class, the requested assertion has the form
\[
\forall G\in\mathcal R\;\forall g,h\in G,\qquad gh=hg.
\]
Here \(\mathcal R\) is an unresolved placeholder: the model-theoretic regularity condition, ambient structures, parameter conventions, and any additional hypotheses are not supplied. Consequently, this part cannot be made fully precise from the record alone.
\end{enumerate}

The source presents these as two questions in the same discussion; it supplies no assertion that they are equivalent or that either implies the other.
\subsection{Short English statement}
The first question asks whether every first-order definable subset of the complex exponential field, allowing complex parameters, is countable or has countable complement; the record attributes this conjecture to Zilber. The second asks whether all regular groups are commutative, but the intended definition of regularity is missing.
\subsection{Sources}
\begin{itemize}
\item Local UnsolvedMath registry: supplied complete record for permanent ID \(30002261\), source problem code \texttt{catalogue entry}, title \emph{Quasiminimal Complex Exponentiation and Abelian Regular Groups}. The record gives proposed year \(2013\), attributes the first conjecture to Zilber, and names no proposer for the combined entry.
\item \emph{Model Theory: Groups, Geometry, and Combinatorics} (2013), Oberwolfach Reports, pp.\ 5--66, DOI \texttt{10.4171/owr/2013/01}: \url{https://doi.org/10.4171/owr/2013/01}. Also listed among the supplied literature sources as \emph{Model Theory, Oberwolfach Report 01/2013}. The registry labels the statement \texttt{corrected\_verified} and says that the official report poses both questions in the same model-theory discussion. These are supplied verification claims, not independently verified here.
\item \emph{Around Podewski's Conjecture}, supplied literature reference: \url{https://doi.org/10.4064/fm222-2-4}.
\item Supplied literature assessment, checked on \texttt{2026-08-21}: the registry labels the entry \texttt{partially\_solved} and states that minimal groups are abelian, that substantial results hold for quasi-minimal and generically stable settings, and that the assertion that every regular group is abelian remains open. This assessment is supplied dated curation derived from the OWR material, not an independently verified literature assessment. The supplied background explicitly requests independent verification of time-sensitive status and cited claims.
\end{itemize}
\end{document}
